\documentclass[11pt]{article}

\usepackage[margin=1in]{geometry}
\usepackage{amsmath}
\usepackage{booktabs}
\usepackage{enumitem}
\usepackage{graphicx}
\usepackage{siunitx}
\usepackage{hyperref}
\usepackage{gensymb}


\title{Basic Rolling Test}
\author{Rift / Team 09}
\date{\today}

\begin{document}


\maketitle

\section*{Test ID}
\textbf{T-01BR Basic Rolling Test} 

\section*{Test Objective}
Determine if the rover is capable of a full rolling cycle, which is defined as 3 full rolling motions. Determine if the rover is capable of returning to its initial position after rolling motions.

\section*{Robot Configuration}
\begin{itemize}[noitemsep]
    \item Robot configurations: Standing
    \item Tube pressures: 15 psi
    \item Control mode: Open-loop. Velocity commands are broadcast to the rollers, which locally use a PID controller to achieve the target speed
    \item Electronics status: Batteries are fully charged (above 4V/cell, 12 V for the full battery)
\end{itemize}

\section*{Materials}
\begin{itemize}[noitemsep]
    \item N/A
    
\end{itemize}

\section*{Environment and Setup}
\begin{itemize}[noitemsep]
    \item Surface type: Finished Smooth Concrete
    \item Surface angle: $\psi = 0 {\degree}$ 
    \item External load: None 
    \item Measurement method: Observation of robot performance
\end{itemize}

\section*{Initial Conditions}
\begin{itemize}[noitemsep]
    \item Initial position: $x_0 = 0$
    \item Initial orientation: $\theta_z = 0 {\degree}$
    \item Initial shape state: Standing
\end{itemize}

\section*{Command Sequence}
\begin{itemize}[noitemsep]
    \item Command type: \underline{\hspace{5 cm}} 
    \item Command values: \underline{\hspace{5 cm}} Command values are generated from the simulation, with potential tuning to account for differences in the simulation and the physical robot
    \item Update rate: 3~Hz
    \item Command Steps: 8~s
\end{itemize}

\section*{Measured Quantities}
\begin{itemize}[noitemsep]
    \item Yes/No of rover's capability to perform 3 consecutive rolling motions
    \item Time required to complete the full rolling cycle
    \item Distance of each node from its initial position after rolling motion
\end{itemize}

\section*{Success Criteria (Per Trial)}
\begin{itemize}[noitemsep]
    \item The rover is capable of performing 3 consecutive rolling motions
    \item The rover is capable of returning to its initial position
\end{itemize}

\section*{Test Procedure}
\begin{enumerate}[noitemsep]
    \item Place robot in standing position
    \item Complete pretest checklist
    \item Send command to rover to perform 1 rolling motion
    \item Repeat Step 3 two more times 
    \item Record if the rover successfully completed all 3 motions
    \item Command rover to return to initial position
    \item Record if rover successfully returned to initial position

\end{enumerate}

\section*{Number of Trials}
\begin{itemize}[noitemsep]
    \item Total trials: 3
    \item Time between trials: N/A
\end{itemize}

\section*{Results}
Summarize the observed performance metrics.

\begin{center}
\begin{tabular}{lc}
\toprule
Metric & Pass/Fail\\
\midrule
Ability to perform 3 rolling motions&  \\
 Ability to return to initial position &\\
\end{tabular}
\end{center}

\section*{Observed Failure Modes}
Describe any observed failure behaviors.

\begin{itemize}[noitemsep]
    \item Truss buckles during motion
    \item Rover is not capable of completing rolling motion, center of mass does not move far enough to induce rolling
    \item Spherical joints reach the limit of their range of motion. 
    \item Rover is not capable of returning to initial position
    \item Rover is damaged during rolling motion
    \item Rover reaches singularity during rolling motion
\end{itemize}

\section*{Conclusions}
Provide a concise interpretation of the results.

\textit{Example:} The robot achieved consistent forward motion on flat ground with an average speed of \underline{\hspace{1 cm}}~cm/s. 

\section*{Notes and Limitations}
Document sources of variability or limitations.

\end{document}